\section{The Fire and the Resurrection of the Body}

Christ warns of punishment as well as exclusion. The fire, the weeping, the gnashing of teeth, and the worm that does not die describe the affliction of the condemned under divine judgment. They suffer the loss of blessedness and the punishment pronounced by the Judge. A person considering his own sin must therefore attend to more than the happiness he has failed to seek. He has offended God, and he is answerable for that offense. The mercy offered now delivers him from a just sentence which Christ describes in these terrible terms.

The Fourth Lateran Council confesses that all will rise in their own bodies and receive according to their deeds: the reprobate perpetual punishment with the devil, the elect everlasting glory with Christ (\emph{Firmiter credimus}, DS 801). The contrast is between two ends of the same bodily resurrection, not between persons who rise and persons who disappear. Christ likewise teaches a resurrection of life and a resurrection of judgment (Jn 5:28--29). The body which shares the labor and choices of earthly life will share the person's final condition.

St.~Cyril of Jerusalem unfolds this teaching while instructing those preparing for baptism. In his eighteenth catechetical lecture, he insists upon the identity of the body raised: the very body which is now mortal will be raised, though transformed. The Lord who made it can bring about what exceeds our present experience. Cyril distinguishes the glory of the righteous from the enduring punishment of the sinner; the latter's body endures suffering without being consumed (\emph{Catechetical Lecture} 18.18--19).

Cyril explains why bodily conduct belongs to this account. The mouth can pray or blaspheme. The body can preserve chastity or commit fornication. The hand can steal or give alms. The same body has served the person in these actions and will share in their final outcome. His examples identify powers which can be used either well or badly. Bodily resurrection therefore bears directly on the catechumen's preparation for a changed life: the body that God will raise is already involved in the obedience or disobedience for which the person must answer.

His next exhortation is to care for the body and not misuse it as though it were not our own (18.20). The Creator witnesses what no human observer sees. Secrecy changes the social circumstances of an act; it does not remove the act from the life God knows. A private cruelty, a concealed theft, or an impure use of another person remains the deed of a bodily creature whose whole existence belongs to the Lord. Reverence for the body includes reverence for its destiny and for the bodies of other people.

Cyril speaks to those approaching baptism of the cleansing God gives, and then urges vigilance against future wounds. The body which needs cleansing is the body God receives into the life of grace. His exhortation to guard it follows from that gift. A baptized Christian returning through Penance can apply the same concern to an established bodily habit. Repeated actions have made one response familiar and easy; repentance may require patient changes before temptation becomes acute. God's forgiveness restores the penitent, who must then cease using his body to renew the fault he has confessed.

St.~Augustine considers the scriptural imagery in the twenty-first book of the \emph{City of God}. He receives Christ's repeated warning about unquenchable fire and affirms that the body itself will suffer by fire. He records different interpretations of the worm: some understand bodily affliction, while others understand the anguish of the soul (XXI.9). The interpretations differ over what that particular image describes. Both concern the enduring misery of the condemned.

In the following chapter, Augustine asks how material fire can afflict spirits. He affirms that it can do so by God's power while acknowledging a mode of connection beyond our comprehension (XXI.10). Ordinary experience of fire cannot explain how a spirit suffers from it. Augustine reasons from the Creator's power over both spiritual and bodily creatures. He can affirm the punishment Christ reveals while admitting that he cannot describe the connection which makes it possible.

Augustine also distinguishes the rich man's experience in Luke 16, before the final resurrection, from the bodily punishment at the end. The parable uses eyes, tongue, flame, and a finger carrying water to speak of an experience occurring while the dead remain without their risen bodies. Its reality does not depend upon turning every feature into a diagram.

The fallen angels are punished too. Christ calls the fire prepared for the devil and his angels; they are subjects of divine judgment. Their rebellion ends in punishment under the Creator whose authority they rejected. The sinner condemned with them shares that punishment. In Matthew's sentence, the Judge commands departure into the fire he has prepared, and those he condemns must undergo his sentence. The spirits who tempted people away from obedience also suffer the final defeat of their own disobedience.

The resurrection requires us to consider what we are doing with the body that will be raised. Cyril's examples make amendment possible to identify. A person who has used his hand to steal must cease taking what belongs to another and restore what he can. The same hand can then serve another's need. These are successive acts of the person God calls to salvation, who remains responsible for his bodily life. Care for the body includes directing its powers toward holiness while grace and time permit amendment.

